\documentclass[12pt]{article}
\usepackage[utf8]{inputenc}
\usepackage[T1]{fontenc}
\usepackage[margin=1in]{geometry}
\usepackage{setspace}
\usepackage{xcolor}
\usepackage[colorlinks=true, urlcolor=blue!60!black, linkcolor=blue!60!black]{hyperref}
\setlength{\parindent}{0em}
\setlength{\parskip}{0.6em}
\renewcommand{\baselinestretch}{1.0}
\begin{document}
\begin{center}
{\LARGE \textit{Chisholm v. Georgia}, 2 U.S. (2 Dall.) 419 (1793)\par}
\vspace{0.5em}
{\large Case Brief --- Two-Page Edition\par}
\end{center}

\section{Memory Jogger}
Article III let a citizen of one State sue another State without its consent; the Eleventh Amendment promptly overturned the ruling (\href{https://supreme.justia.com/cases/federal/us/2/419/}{2 U.S. 419}).

\section{Detailed Case Facts}
In 1777 Georgia's commissioners bought war supplies from South Carolina merchant Robert Farquhar and never paid. Alexander Chisholm, Farquhar's executor, was refused payment by the Georgia legislature in 1789.

\section{Procedural History}
A federal circuit court dismissed Chisholm's first suit for lack of jurisdiction. He then sued Georgia in \emph{assumpsit} directly in the Supreme Court; Georgia refused to appear, and only Attorney General Edmund Randolph argued.

\section{Judicial Votes}
4--1 for the plaintiff in seriatim opinions by Chief Justice Jay and Justices Cushing, Wilson, and Blair; Justice Iredell dissented.

\section{Holding}
Article III's grant of jurisdiction over controversies ``between a State and citizens of another State'' covers States as defendants. Because sovereignty lies with the people, Georgia had no immunity, and default judgment would enter unless it appeared.

\section{Analysis of Opinions}
Blair and Cushing read the text literally, rejecting assurances by Hamilton, Madison, and Marshall that States would be only plaintiffs. Wilson and Jay grounded the result in popular sovereignty: here ``there are citizens, but no subjects,'' and a State, like a merchant, must answer for its contracts. Jay noted that Georgia was suing as a plaintiff elsewhere. Iredell dissented because the Judiciary Act authorized only writs known to the common law, under which no compulsory suit lay against the sovereign, and because each State remains ``completely sovereign'' as to its reserved powers. The Eleventh Amendment, ratified in 1795, overturned the decision.

\section{Examples: Future Applications}
\textbf{Same-Side:} a State may still sue an out-of-state supplier as plaintiff, the asymmetry Jay criticized. \textbf{Opposite-Side:} a supplier's federal damages suit against Louisiana for an unpaid contract is barred by the Eleventh Amendment and \emph{Hans}. \textbf{Fence-Sitter:} whether a separately incorporated hospital authority whose larger judgments the State pays is an arm of the State under \emph{Galette}.

\section{Critique}
The majority read Article III literally, yet modern sovereign immunity rests on unwritten structural postulates drawn from Iredell's dissent. Critics also argue that the Eleventh Amendment merely repealed Citizen--State diversity jurisdiction, leaving federal-question suits intact.

\section{Key Quotations}
\begin{quote}
``No such ideas obtain here; at the Revolution, the sovereignty devolved on the people, and they are truly the sovereigns of the country, but they are sovereigns without subjects (unless the African slaves among us may be so called), and have none to govern but themselves; the citizens of America are equal as fellow citizens, and as joint tenants in the sovereignty.''
\end{quote}
\noindent --- Chief Justice Jay, 2 U.S. at 471--472.

\end{document}
